\documentclass{article}

\usepackage{amssymb,amsmath}

\begin{document}

General model assumptions:
\begin{itemize}
    \item the fluid density and viscosity are constant in time and space (implies incompressibility);
    \item the fluid flow is laminar;
    \item the column is operated under constant conditions (e.g., temperature, flow rate);
    \item the process is isothermal (i.e., there are no thermal effects);
    \item diffusion does not depend on the concentration of the components, viscosity of the fluid, or pressure;
    \item the fluid in the particles is stagnant (i.e., there is no convective flow);
    \item the porous particles are spherical, rigid, and do not move;
    \item the particles are homogeneous and of uniform porosity;
    \item all components access the same particle volume;
    \item the partial molar volumes are the same in mobile and solid phase;
    \item the binding model parameters do not depend on pressure and are constant along the column;
    \item the solvent is not adsorbed;
\end{itemize}

Specific model assumptions:
\begin{itemize}
    \item the particles form a continuum inside the column (i.e., there is interstitial and particle volume at every point in the column);
    \item the column is radially symmetric and homogeneous (i.e., concentration profiles and parameters only depend on the axial position);
    \item transverse concentration variations equilibrate sufficiently rapidly relative to the flow-directional transport for their effects to be represented by an effective flow-directional dispersion coefficient;
    \item the interstitial liquid phase concentration is spatially constant on the outer particle surface;
    \item there is no surface diffusion (i.e. adsorbed molecules are spatially constant);
\end{itemize}

SMA binding model assumptions:
\begin{itemize}
    \item charged species interact with charged ligands on the resin;
    \item electroneutrality is enforced via displacing counter-ions;
    \item steric sheilding of binding sites may take place;
\end{itemize}

\section*{General Rate Model with SMA binding}
Consider a cylindrical column of length $L^{\mathrm{b}} > 0$ packed with spherical particles, and observed over a time interval $(0, T^{\mathrm{end}})$.
In the interstitial volume, mass transfer is governed by the following convection-diffusion-reaction equations in $(0, T^{\mathrm{end}})\times (0, L^{\mathrm{b}})$ and for all components $i\in\{1, \dots, N^{\mathrm{c}} \}$
\begin{align}
    \varepsilon^{\mathrm{b}} \frac{\partial c^{\mathrm{b}}_{i}}{\partial t}
    &= - u \frac{\partial \left( \varepsilon^{\mathrm{b}} c^{\mathrm{b}}_{i} \right)}{\partial z}
        + \frac{\partial}{\partial z} \left( \varepsilon^{\mathrm{b}} D^{\mathrm{ax}}_{i} \frac{\partial c^{\mathrm{b}}_{i}}{\partial z} \right)
        - \left(1 - \varepsilon^{\mathrm{b}} \right) \frac{3}{R^{\mathrm{p}}} k^{\mathrm{f}}_{i} \left(c^{\mathrm{b}}_{i} - \left. c^{\mathrm{p}}_{i} \right|_{r = R^{\mathrm{p}}} \right),
\end{align}
with boundary conditions
\begin{alignat}{2}
    u c^{\mathrm{in}}_{i}
    &= \left.\left( u c^{\mathrm{b}}_{i} - D^{\mathrm{ax}}_{i} \frac{\partial c^{\mathrm{b}}_{i}}{\partial z}\right)\right|_{z=0}
    & & \qquad\text{on }(0, T^{\mathrm{end}}),
    \\
    0
    &= - D^{\mathrm{ax}}_{i} \left. \frac{\partial c^{\mathrm{b}}_{i}}{\partial z} \right|_{z=L^{\mathrm{b}}}
    & & \qquad\text{on }(0, T^{\mathrm{end}}).
\end{alignat}
Here, $t\in (0, T^{\mathrm{end}})$ is the time coordinate, $z\in (0, L^{\mathrm{b}})$ is the axial cylinder coordinate, $R^{\mathrm{p}}> 0$ is the particle radius, $c^{\mathrm{b}}_{i}\colon (0, T^{\mathrm{end}})\times (0, L^{\mathrm{b}}) \mapsto \mathbb{R}$ is the bulk liquid concentration, $\varepsilon^{\mathrm{b}}\in (0, 1)$ is the bulk porosity, $u> 0$ is the interstitial velocity, $D^{\mathrm{ax}}_{i}\geq 0$ is the axial dispersion coefficient, and $k^{\mathrm{f}}_{i}\geq 0$ is the film diffusion coefficient.
In the particles, mass transfer is governed by diffusion-reaction equations in $(0, T^{\mathrm{end}}) \times (0, L^{\mathrm{b}})\times (0, R^{\mathrm{p}})$ and for all non-salt components
\begin{align}
    \frac{\partial c^{\mathrm{p}}_{i}}{\partial t}
    &= \frac{1}{r^{2}} \frac{\partial }{\partial r} \left( r^{2} D^{\mathrm{p}}_{i} \frac{\partial c^{\mathrm{p}}_{i}}{\partial r} \right)
        - \frac{1 - \varepsilon^{\mathrm{p}}}{\varepsilon^{\mathrm{p}}}\left(k^{\mathrm{a}}_{i} c^{\mathrm{p}}_{i} \left( \frac{\bar{q}_{0}}{q^{\mathrm{ref}}} \right)^{\nu_{i}} - k^{\mathrm{d}}_{i} c^{\mathrm{s}}_{i} \left( \frac{c^{\mathrm{p}}_{0}}{c^{\mathrm{ref}}} \right)^{\nu_{i}} \right),
    \\
    \frac{\partial c^{\mathrm{s}}_{i}}{\partial t}
    &= k^{\mathrm{a}}_{i} c^{\mathrm{p}}_{i} \left( \frac{\bar{q}_{0}}{q^{\mathrm{ref}}} \right)^{\nu_{i}}
        - k^{\mathrm{d}}_{i} c^{\mathrm{s}}_{i} \left( \frac{c^{\mathrm{p}}_{0}}{c^{\mathrm{ref}}} \right)^{\nu_{i}},
\end{align}
with boundary conditions
\begin{align}
    - \left. \left( D^{\mathrm{p}}_{i} \frac{\partial c^{\mathrm{p}}_{i}}{\partial r} \right) \right|_{r=0}
    &= 0,
    \\
    \varepsilon^{\mathrm{p}} \left. \left( D^{\mathrm{p}}_{i} \frac{\partial c^{\mathrm{p}}_{i}}{\partial r} \right)\right|_{r = R^{\mathrm{p}}}
    &= k^{\mathrm{f}}_{i} \left. \left( c^{\mathrm{b}}_{i} - c^{\mathrm{p}}_{i} \right|_{r = R^{\mathrm{p}}} \right).
\end{align}
Further details on the binding model can be found in the original formulation Brooks and Cramer (1992), extension to multiple bound states in Diedrich et al. (2017).
Here, $r$ is the radial particle coordinate, $c^{\mathrm{p}}_{i}\colon (0, T^{\mathrm{end}}) \times (0, L^{\mathrm{b}})\times (0, R^{\mathrm{p}}) \mapsto \mathbb{R}$ is the particle liquid concentration, $c^{\mathrm{s}}_{i}\colon (0, T^{\mathrm{end}}) \times (0, L^{\mathrm{b}})\times (0, R^{\mathrm{p}}) \mapsto \mathbb{R}$ is the particle solid concentration, $\varepsilon^{\mathrm{p}}\in (0, 1)$ is the particle porosity, $D^{\mathrm{p}}_{i}> 0$ is the particle diffusion coefficient, $k^{\mathrm{a}}_{i}$ is the adsorption rate constant, $k^{\mathrm{d}}_{i}$ is the desorption rate constant, $\nu_{i}$ is the characteristic charge, $\sigma_{i}$ is the steric factor, $\Lambda$ is the ionic capacity (binding site concentration), $q^{\mathrm{ref}}$ is the reference solid phase concentration, and $c^{\mathrm{ref}}$ is the reference liquid phase concentration.
The number of available binding sites $\bar{q}_{0}$ is given by
\begin{align}
    \bar{q}_{0}
    &= \Lambda
        - \sum_{n=1}^{N^{\mathrm{c}} - 1} \left( \nu_{n} + \sigma_{n} \right) c^{\mathrm{s}}_{n}.
\end{align}
For the salt component, mass transfer is governed by diffusion-reaction equations in $(0, T^{\mathrm{end}}) \times (0, L^{\mathrm{b}})\times (0, R^{\mathrm{p}})$
\begin{align}
    \frac{\partial c^{\mathrm{p}}_{i}}{\partial t}
        + \frac{1 - \varepsilon^{\mathrm{p}}}{\varepsilon^{\mathrm{p}}} \frac{\partial c^{\mathrm{s}}_{i}}{\partial t}
    &= \frac{1}{r^{2}} \frac{\partial }{\partial r} \left( r^{2} D^{\mathrm{p}}_{i} \frac{\partial c^{\mathrm{p}}_{i}}{\partial r} \right),
    \\
    0
    &= - c^{\mathrm{s}}_{0}
        + \Lambda
        - \sum_{n=1}^{N^{\mathrm{c}} - 1} \nu_{n} c^{\mathrm{s}}_{n},
\end{align}
where the counter-ion concentration $c^{\mathrm{s}}_{0}$ satisfies the electroneutrality constraint.
Boundary conditions are
\begin{align}
    - \left. \left( D^{\mathrm{p}}_{i} \frac{\partial c^{\mathrm{p}}_{i}}{\partial r} \right) \right|_{r=0}
    &= 0,
    \\
    \varepsilon^{\mathrm{p}} \left. \left( D^{\mathrm{p}}_{i} \frac{\partial c^{\mathrm{p}}_{i}}{\partial r} \right)\right|_{r = R^{\mathrm{p}}}
    &= k^{\mathrm{f}}_{i} \left. \left( c^{\mathrm{b}}_{i} - c^{\mathrm{p}}_{i} \right|_{r = R^{\mathrm{p}}} \right).
\end{align}
Consistent initial values for all solution variables (concentrations) are defined at $t = 0$.
This model is a particular instance of the unified mathematical framework introduced in Breuer et al. (2026), which provides a rigorous justification of its assumptions and a systematic derivation of its governing equations.
\begin{thebibliography}{99}
    \bibitem{brooks1992steric} Brooks, Clayton A and Cramer, Steven M (1992). Steric mass-action ion exchange: displacement profiles and induced salt gradients. AIChE Journal. DOI: https://doi.org/10.1002/aic.690381212.
    \bibitem{DIEDRICH201760} Juliane Diedrich and William Heymann and Samuel Leweke and Stephen Hunt and Robert Todd and Christian Kunert and Will Johnson and Eric von Lieres (2017). Multi-state steric mass action model and case study on complex high loading behavior of mAb on ion exchange tentacle resin. Journal of Chromatography A. DOI: https://doi.org/10.1016/j.chroma.2017.09.039.
    \bibitem{breuer2026unified} Jan Michael Breuer and Samuel Leweke and Johannes Schmölder and Antonia Berger and Gregor J. Gassner and Eric von Lieres (2026). A unified and composable framework for modeling packed-bed liquid chromatography. arXiv preprint arXiv:2610.33269. DOI: https://doi.org/10.48550/arXiv.2610.33269.
\end{thebibliography}

\end{document}
